\section{Expériences Professionnelles \& Associatives}
\vspace{0.1cm}

\subsection{{JP Morgan}}
\vspace{-0.2cm}
\begin{leftbar}
    \subtext{\textbf{AI Research Associate \hfill Février 2025 - Présent}}
    \begin{zitemize}
        \item Chercheur au sein de l'équipe AI Research
    \end{zitemize}
    \subtext{\textbf{AI Research Scientist --- stagiaire \hfill Avril 2024 - Septembre 2024}}
    \begin{zitemize}
        \item Publication à MATH-AI 2024 (conférence de NeurIPS) --- outstanding paper award
        \item Travail sur l'utilisation de LLMs (Large Language Models) pour des tâches de raisonnement mathématique
        \item Développement d'un système d'algèbre symbolique pour la génération de données synthétiques en langage naturel
        \item Développement d'un système de méta-programmation en Lean pour la génération de données synthétiques
    \end{zitemize}
    \vspace{-0.2cm}
\end{leftbar}
\vspace{-0.2cm}

\subsection{{EDF}}
\vspace{-0.2cm}
\begin{leftbar}
    \subtext{\textbf{Data Scientist --- stagiaire \hfill Janvier 2023 - Juillet 2023}}
    \begin{zitemize}
        \item Chargé de la conception et de l'implémentation d'un outil de détection d'anomalies dans des séries temporelles
        \item Conception d'un banc de test de méthodes de détection d'anomalies et réalisation d'un comparatif
        \item Développement d'une API exposant des méthodes de détection d'anomalie avec stockage en base
        \item Développement d'une interface web pour la visualisation des anomalies détectées
    \end{zitemize}
    \vspace{-0.2cm}
\end{leftbar}
\vspace{-0.2cm}

\subsection{{Theodo --- groupe M33}}
\vspace{-0.2cm}
\begin{leftbar}
    \subtext{\textbf{Software Engineer --- stagiaire \hfill Juin 2022 - Décembre 2022}}
    \begin{zitemize}
        \item Développeur full stack \textit{(AWS Serverless / Symphony / Spring Boot --- React / Angular)}
        \item Développement dans un environnement agile
        \item Développement d'un \textit{back-office} pour un acteur du e-commerce: une plateforme pour intégrer de nouveaux produits au site avec un système de reconaissance des produits déjà existants
        \item Développement d'une plateforme d'accompagnement des entrepreneurs pour une banque d'investissement
    \end{zitemize}
    \vspace{-0.2cm}
\end{leftbar}
\vspace{-0.2cm}

\subsection{{Club informatique de l’École nationale des Ponts et Chaussées}}
\vspace{-0.2cm}
\begin{leftbar}
    \subtext{\textbf{Responsable technique --- bénévole \hfill Avril 2021 - Juin 2022}}
    \begin{zitemize}
        \item Développement d'applications webs pour l'école et les élèves
        \item Administration et maintenance des serveurs et du \textit{cloud}, déploiement des sites créés par l'association et par des étudiants
        \item Réparation de matériel informatique pour les élèves, remise en état de matériel informatique pour des réfugiés
    \end{zitemize}
    \vspace{-0.2cm}
\end{leftbar}
\vspace{-0.2cm}
